\documentclass[11pt]{article}
\usepackage{amsmath,amssymb,booktabs}
\setlength{\textwidth}{6.4in}
\setlength{\oddsidemargin}{0.05in}
\setlength{\evensidemargin}{0.05in}
\setlength{\textheight}{9in}
\setlength{\topmargin}{-0.4in}
\newcommand{\tr}{\operatorname{tr}}
\newcommand{\KL}{\operatorname{KL}}
\newcommand{\NLL}{\operatorname{NLL}}
\newcommand{\norm}[1]{\left\lVert #1\right\rVert}
\title{JLZ v11: Joint Local Increments with Native Ridge-Cost Allocation}
\author{Method and experiment contract, revision 1}
\date{October 4, 2026}
\begin{document}
\maketitle

\paragraph{Scope.}
The method replaces MEMIT-H's single upper-layer target and remaining-layer
divisor with jointly optimized local increments at every configured write layer.
A ridge-derived cost supplies an allocation signal. It preserves native input
sentences, subject lookup, task readouts, current-to-entry KL direction, the ridge
writer, causal key refresh at commit, and history accumulation. It does not claim
to solve every source of semantic locality loss. All configured layers remain
eligible; concentration on a small subset is an allowed outcome, not a gate.

\section{Entry state and local variables}
At sequential batch $t$, let $W_t$ be the entry model and $H_{t,l}$ the previous
key-history covariance. Let $B$ be the actual number of requests, including a
partial final batch; $\mathcal L$ the configured write layers; $m=|\mathcal L|$;
and $d_l$ the layer output dimension. All of these are runtime quantities.
For request $r$, let $h^0_{lr}$ be the canonical entry subject's full
transformer-block output and $a_{lr}=\norm{h^0_{lr}}_2>0$. The model, anchors,
entry KL teacher $p^0_{K,r}$, and history are fixed during fitting.

Inject the same trainable $\delta_{lr}\in\mathbb R^{d_l}$ at the subject-last
lookup of every native rewrite and KL context of request $r$. Define
$D_l=[\delta_{l1},\ldots,\delta_{lB}]$. Each request's forward contains all its
layer increments simultaneously; gradients propagate through the effects of
lower-layer injections on upper-layer states. There is no separate task-success
objective for each layer. Model weights are not optimizer variables.

The injection site is the full block output; the writer key is the down-projection
input. A model adapter must verify their local additive correspondence.
Token IDs, target positions, padding and position semantics, NLL readout, final
KL readout, and native context definitions are preserved. Evaluation paraphrases
and neighborhood examples never enter fitting or coefficient selection.

\section{Task objective}
The virtual subject-injection loss is
\begin{equation}
 L_N(D)=\frac1B\sum_{r=1}^{B}\left[
 \NLL_r^v+\lambda_K\KL(p^v_{K,r}\Vert p^0_{K,r})
 +\lambda_n\sum_{l\in\mathcal L}
       \frac{\norm{\delta_{lr}}_2}{a_{lr}^2}\right].
 \label{eq:native}
\end{equation}
NLL uses the native target-token mean followed by its context mean. KL uses the
full vocabulary and the current-to-entry direction. Decay is the unsquared
Euclidean norm divided by the squared canonical entry anchor. The layerwise
sum extends the native single-increment penalty; it does not imply identical
effective regularization or an identical total edit budget.

\section{Frozen entry geometry and full-batch allocation cost}
Extract entry rewrite keys and take the native nested context mean to obtain
$\widehat K_l\in\mathbb R^{k_l\times B}$. KL rows do not enter this mean.
In the current six-rewrite-context profile, key weights are
$(0.5,0.1,0.1,0.1,0.1,0.1)$, while rewrite NLL weights are each $1/6$.
These are distinct reductions. Define
\begin{align}
 A_l&=\lambda_C C_{0,l}+H_{t,l},\\
 \widehat P_l&=(A_l+\widehat K_l\widehat K_l^\top)^{-1}\widehat K_l,
 &\widehat M_l&=\widehat P_l^\top\widehat K_l.
\end{align}
Use solves, not explicit inverses. This geometry is constructed once at batch
entry and frozen across candidates; it is rebuilt for the next sequential batch.

For symmetric positive-definite $A_l$, the exact optimum at this fixed geometry is
\begin{align}
 \widehat C_l(D_l)
 &=\min_U\left\{\norm{U\widehat K_l-D_l}_F^2+
                         \tr(U A_lU^\top)\right\}\\
 &=\tr\left[D_l(I-\widehat M_l)D_l^\top\right].
 \label{eq:ridgecost}
\end{align}
In particular, with
$G_l=\widehat P_l^\top A_l\widehat P_l$ and
$E_l=(\widehat M_l-I)(\widehat M_l-I)^\top$,
$G_l+E_l=I-\widehat M_l$. The method retains this combined ridge objective,
not the preservation energy $G_l$ alone. For a stable norm representation, set
$S_l=\widehat K_l^\top A_l^{-1}\widehat K_l$ and factor
$I+S_l=T_lT_l^\top$. Then
\begin{equation}
 \widehat C_l(D_l)=\norm{T_l^{-1}D_l^\top}_F^2.
\end{equation}
Native covariance/history rounding can produce slight asymmetry. Qualification
must check agreement with native asset solves; arbitrary symmetrization or
unreported jitter is not permitted. A failed tolerance is a technical failure.

The main method uses the A-form sum of roots:
\begin{align}
 \sigma_l^2&=\frac1B\sum_r a_{lr}^2,&
 c_l(D_l)&=\frac{\sqrt{\widehat C_l(D_l)}}{\sqrt B\,\sigma_l},\\
 \mathcal R_A(D)&=\sum_{l\in\mathcal L}c_l(D_l),&
 J(D)&=L_N(D)+\lambda_{\rm alloc}\mathcal R_A(D).
 \label{eq:joint}
\end{align}
The minimum-norm subgradient at a zero norm is zero; no smoothing term is added.
There is no additional root-of-sum-of-squares arm. Full-batch off-diagonal terms
are retained, so the allocation cost couples requests even though native task
losses can be computed separately. A microbatch-specific ridge or diagonal
approximation would define a different method.

\paragraph{Exact identity versus proxy.}
Equation~\eqref{eq:ridgecost} is an exact ridge identity for the frozen entry
geometry. Its use during fitting is a proxy for the eventual commit cost:
upper-layer keys change after lower-layer writes. The method does not label
this proxy as the cost of the actual deployed update.

\section{Optimizer, projection, and stopping}
Reparameterize
\begin{equation}
 \delta_{lr}=s_{lr}q_{lr},\qquad
 s_{lr}=\frac{a_{lr}}{\sqrt{d_lm}},\qquad q_{lr}^{(0)}=0.
\end{equation}
Use Adam on $q$: learning rate $0.1$, betas $(0.9,0.999)$,
$\epsilon=10^{-8}$, no weight decay, no warm-up, and no foreach implementation.
The default model, $q$, $D$, and moments are FP32; geometry and allocation
reductions are FP64. This reparameterization is an explicit adaptation, not
the same physical step as native Adam on $\delta$.

Log the request-mean objective $J$. Pass a single request-sum gradient to Adam,
as in the existing JLZ convention:
\begin{equation}
 (q_{lr}.\mathrm{grad})=B\,s_{lr}\frac{\partial J}{\partial\delta_{lr}}.
\end{equation}
Apply $B$ and $s_{lr}$ exactly once. Accumulate native task gradients over
microbatches and add the full-batch allocation gradient once per candidate.

After every Adam update, project each physical increment onto
\begin{equation}
 \norm{\delta_{lr}}_2\leq c_{\rm native}a_{lr},
 \qquad c_{\rm native}=0.75\quad\hbox{in the current profile},
\end{equation}
and restore $q_{lr}=\delta_{lr}/s_{lr}$ while retaining Adam moments.
There is no additional total-layer budget or clipping of model hidden states.

There are at most 25 candidate evaluations and 24 optimizer updates. Candidate
1 is $q=0$. At each candidate, first evaluate the full $J$ and reject nonfinite
or corrupted states as technical failures. Stop the whole logical batch when
$J_{\rm mean}<0.05$, including allocation cost, or after candidate 25. Otherwise
backpropagate, update Adam, project, and evaluate the next candidate. Adopt the
last evaluated finite candidate even if it exhausts the budget. Do not select
an earlier candidate by evaluation performance or add steps after the budget.

This is a joint extension of native total-loss stopping, not native independent
request stopping. There is no finished-request freeze, NLL-only sufficiency
stop, or ACC/PS/NS-based gate. The no-allocation control applies the same rule
to its own objective with $\lambda_{\rm alloc}=0$.

\section{Sequential native ridge commit}
Freeze the fitted $D_l$. Inside a transaction based on $(W_t,H_t)$, process all
write layers in native causal order:
\begin{enumerate}
 \item Extract rewrite keys from the current model, which already contains
       the committed updates at earlier layers, and form native mean keys $K_l^a$.
 \item Solve $P_l^a=(A_l+K_l^aK_l^{a\top})^{-1}K_l^a$ using the native writer.
 \item Form $U_l=D_lP_l^{a\top}$ and apply the native orientation, cast, and
       addition to the selected weight. Proceed to the next layer in the updated model.
\end{enumerate}
The writer's requested residual is exactly $D_l$. There is no remaining-layer
divisor, absolute virtual target minus actual hidden state, residual top-up,
or inverse-$M$ compensation. Earlier under-realization is not automatically
assigned to a later layer. Entry $\widehat P_l$ must not be reused for commit.

At actual mean keys, direct realization is
\begin{equation}
 Y_l=U_lK_l^a=D_lM_l^a,\qquad M_l^a=P_l^{a\top}K_l^a.
\end{equation}
Planned and realized shares need not agree. If an entire $D_l$ is zero, then
$U_l=0$; an individual zero column does not imply zero effect at that request's
key because other columns can contribute. Report direct local action $U_lk_l$
separately from full hidden differences inherited from lower layers.

After all writes, re-extract native mean keys from the final model and append
$H_{t+1,l}=H_{t,l}+K_{{\rm hist},l}K_{{\rm hist},l}^{\top}$ once in CPU FP32.
Key reuse requires prior final-key parity. Verify preservation of nonselected
weights and rollback. There is no history replay or replay teacher.

\paragraph{Remaining discrepancy.}
Fitting applies a common increment at subject positions; committed weights act
at every token and only partly realize the request at its mean key. Consequently,
local targets reduce the target-location mismatch without guaranteeing exact
realization or fit-to-commit task parity. No actual-writer injection (T$'$),
pulse, reference loss, or all-token teacher is introduced.

\section{Native boundary and efficiency}
Native text, scoring, canonical anchor convention, ridge formula, covariance,
history, and actual key refresh are retained. Explicit extensions are the
multi-layer direct-increment path; summed layer norms and caps; $q$ Adam;
full-batch allocation; joint stopping; and replacing upper-target residual
feedback with fixed local increments. Equal coefficients do not imply equal
effective strength. No B1 coefficient calibration, strength matching, or sweep
is part of this design.

Entry geometry is computed once and the sequential actual writer once after
fitting, removing candidate-by-candidate actual builder solves and weight
materialization. Verified token-preserving padding trims, selected-position
full-vocabulary heads, valid prefix reuse, and early key-extraction exit are
allowed. Logical batch normalization, contexts, vocabulary, dtypes, and candidate
rules must remain unchanged. Report logical candidates, actual forward/backward
counts, time, and memory separately; microbatch utilization gains are not
automatically FLOP savings.

Model/benchmark adapters supply dimensions, layers, native hyperparameters,
context definitions, and injection/readout modules. Unsupported architectures
fail qualification explicitly rather than using silent fallbacks.

\section{Qualification and the initial 2,000-edit experiment}
The method is independent of batch size and model. The initial execution profile
uses Llama-3-8B-Instruct and CounterFact with layers L4--L8, loss layer 31,
$\lambda_K=0.0625$, $\lambda_n=0.5$, $\lambda_C=15000$, and the native
$0.75$ cap. The main allocation coefficient is locked by the accompanying
experiment configuration; it is not selected using evaluation results.

Proceed through CPU math checks, small real-model qualification, a BS2-by-2
pilot, and actual BS100 admission. Technical parity, finite gradients, correct
masking, and state/commit integrity govern admission; accuracy and layer shares
do not. The main comparison uses the same ordered 2,000 requests in 20 logical
batches of 100 for all three methods:
\begin{center}
\begin{tabular}{lll}
\toprule
Method & Allocation term & Comparison meaning\\
\midrule
JLZ v11 main & $\lambda_{\rm alloc}\mathcal R_A$ & Full proposed pipeline\\
JLZ v11 NOALLOC & $0$ & Added writer-cost effect\\
Native MEMIT-H & Native procedure & Baseline\\
\bottomrule
\end{tabular}
\end{center}
NOALLOC does not impose uniform allocation: task losses and norm penalties can
still favor unequal increments. Main versus NOALLOC measures the effect of an
added cost, including both regularization strength and allocation changes.
NOALLOC versus MEMIT-H compares the joint local-target pipeline; it is not an
isolated causal test of target location alone.

Report rewrite and paraphrase ACC alongside preference RS/PS, neighborhood
ACC and NS, their exact denominators and scoring definitions, NLL, immediate
cohort performance, and subsequent retention. Preserve ACC as a separate
outcome; preference saturation does not imply that further optimization is
useless. Include planned and realized allocation, actual-versus-entry cost,
clipping, stopping, and compute telemetry. The authoritative implementation and
experiment documents specify numerical tolerances, schedules, and artifacts.
\end{document}
